\chapter{Section B: Tracing Problems}\label{ch:tracing}

Determine the exact output of each program, including spaces and line breaks. Each one
hides at least one trap; the most common are listed in the checklist below. Unless
stated, the code compiles and runs without error. Programs that rely on a library of
Part~I (\cc{stackqueue.h}, \cc{deque.h}, \cc{sll.h}, \cc{dll.h}, \cc{circular.c}) or on
the solution of a problem of Chapter~\ref{ch:problems} say so, and the behaviour of the
functions they call is stated before the program. The headers \cc{stdio.h} and
\cc{stdlib.h} are included everywhere and not shown.

\para{Tracing checklist} (i)~Evaluation order of function arguments is unspecified; the
programs here use sequential statements to avoid it. (ii)~Pass by value copies a struct
including an embedded array, but a pointer member still points to the same nodes.
(iii)~A loop bound such as \cc{i<=n} or \cc{i<size-1} changes the count by one.
(iv)~\cc{printf} adds nothing: look for trailing spaces and \cc{\textbackslash n}.
(v)~Integer division truncates toward zero, and \cc{\%} keeps the sign of the left
operand. (vi)~A failed \cc{pop} of the course library returns 0 and leaves its output
variable \emph{unchanged}. (vii)~Keep a table with one column per variable and one row
per step; do not trace in your head.

\tracing{1}{Array list, return values and stale cells}
A bounded array list with capacity 5 and 1-based \cc{setItem} and \cc{itemAt}. Give the
exact output of the program.
\programfile{tracing/t01.c}

\tracing{2}{What does the callee really change?}
Three structs are passed by value, by pointer and as a raw array. For each of the lines
\cc{A}--\cc{E} state what is printed.
\programfile{tracing/t02.c}

\tracing{3}{Counting loop iterations}
Give the six counters. Then state each count as a function of $n$ in $\Theta$-notation.
\programfile{tracing/t03.c}

\tracing{4}{The raw array after erase and insert}
The class prints its \emph{whole} array, including the cells beyond \cc{size}. Positions
are 0-based.
\programfile[language=C++]{tracing/t04.cpp}

\tracing{5}{Circular queue: wrap-around}
A circular queue of capacity 4. \cc{show} prints the fields and the raw array.
\programfile{tracing/t05.c}

\tracing{6}{SLL with head, tail and size}
The program uses \cc{sll.h} (Listing~\ref{lst:sll}) and \cc{popFront}/\cc{popBack} of
Listing~\ref{lst:sllpop}: each pop returns 1 and stores the removed item in \cc{*x}, or
returns 0 and leaves \cc{*x} unchanged when the list is empty. \cc{sdisplay} prints the
list as \cc{[3 5]}. \cc{\%-11s} pads the operation name to 11 characters.
\programfile{tracing/t06.c}

\tracing{7}{One star or two?}
\programfile{tracing/t07.c}

\tracing{8}{Erasing from a doubly linked list}
The program uses \cc{dll.h} (Listing~\ref{lst:dll}: \cc{display} prints the items
separated by spaces and a newline; \cc{displayBackward} prints them from the tail) and
\cc{eraseItems(list *l, int x)} of Problem~\ref{prob:13}, which removes \emph{every} node
holding \cc{x}. The departmental output format ends each item with a space.
\programfile{tracing/t08.c}

\tracing{9}{Recursion on a list}
The lists are built with a pointer to the last \cc{next} field. Give the three lines.
\programfile{tracing/t09.c}

\tracing{10}{Stacks and queues in the same function}
The program uses \cc{stack<int>} and \cc{queue<int>}: \cc{top()}/\cc{front()} read an
item without removing it, \cc{pop()} removes it and returns nothing.
\programfile{tracing/t10.cpp}

\tracing{11}{Infix to postfix with a depth counter}
A variant of the shunting-yard algorithm on letters. Give the output string and the
maximum size the operator stack reaches.
\programfile{tracing/t11.cpp}

\tracing{12}{A queue from two stacks}
Give the output, and state how many items are moved from \cc{in} to \cc{out} in total.
\programfile{tracing/t12.cpp}

\tracing{13}{Which operand is popped first?}
Both functions evaluate a postfix expression of single digits with integer division.
They differ only in the order of the two pops. Give the table.
\programfile{tracing/t13.cpp}

\tracing{14}{Stack permutations}
\cc{produce} tries to pop the numbers in the order \cc{want[]} while pushing
$1,2,3,4$ in order. Give the five lines and the count.
\programfile{tracing/t14.cpp}

\tracing{15}{The stock spanner}
The program uses \cc{StockSpanner} of Problem~\ref{prob:37}: \cc{next} pops every
(price, span) pair whose price is $\le$ today's price, adds their spans to 1, and pushes
today's pair. \cc{show} prints the pairs from bottom to top.
\programfile{tracing/t15.cpp}

\tracing{16}{The mid-pointer list}
The program uses \cc{mlist} of Problem~\ref{prob:24}: a singly linked list with
\cc{head}, \cc{mid}, \cc{tail} and \cc{size}, where \cc{middle(l)} returns the item at
index $\lfloor(\cc{size}-1)/2\rfloor$. Give the \cc{mid} column; the list contents follow
from the operations.
\programfile{tracing/t16.c}

\tracing{17}{Binary search: counting the steps}
\cc{lowerBound} returns the index of the first item $\ge x$ on the half-open range
$[\cc{lo},\cc{hi})$.
\programfile{tracing/t17.c}

\tracing{18}{Josephus with a twist}
Six people, every third one leaves. Give the order of removal and the winner, and check
the winner with the recurrence of Chapter~\ref{ch:queues}.
\programfile{tracing/t18.cpp}

\tracing{19}{Splicing lists}
The program uses \cc{sll.h} and the two functions of Problem~\ref{prob:18}:
\cc{spliceAfter(l, p, m)} moves all nodes of \cc{m} after node \cc{p} of \cc{l}, and
\cc{splice(l, pos, m)} moves them so that the first one becomes item \cc{pos} of \cc{l}
($0\le\cc{pos}\le\cc{l->size}$; it returns 0 and does nothing for an invalid
\cc{pos}). Both leave \cc{m} empty.
\programfile{tracing/t19.c}

\tracing{20}{Shallow and deep copies}
\programfile{tracing/t20.c}

\tracing{21}{Circular deque: the raw array}
A circular deque of capacity 5. \cc{pushFront} steps the front back by adding 4 instead
of subtracting 1. The last line prints three C expressions.
\programfile{tracing/t21.cpp}

\tracing{22}{ArrayDeque: which side moves?}
Positional insertion in a circular array of capacity 8 that moves the shorter side.
\cc{add} returns the number of items moved.
\programfile{tracing/t22.cpp}

\tracing{23}{Sliding window maximum}
The program uses \cc{deque.h} (Listing~\ref{lst:deque}) and prints the deque of indices
after each step of the algorithm of Problem~\ref{prob:53}. Give the eight lines. Watch the
\emph{two} equal values.
\programfile{tracing/t23.cpp}

\tracing{24}{A circular list as a queue}
The program uses Listing~\ref{lst:circular}: \cc{enq(\&last, x)} links a new node after
\cc{last} and makes it the new \cc{last}; \cc{deq(\&last, \&x)} removes \cc{last->next}
(the front), stores its item in \cc{x} and returns 1, or returns 0 when the list is empty.
The main program also moves the pointer by hand.
\programfile{tracing/t24.c}

\tracing{25}{The sentinel deque, and a shallow copy}
The program uses the linked deque of Problem~\ref{prob:51} (Listing~\ref{lst:linkeddeque}).
Print 1 for true and 0 for false in the comparisons.
\programfile{tracing/t25.cpp}
